\documentclass[10pt,letterpaper]{article}
\usepackage[T1]{fontenc}
\usepackage{amssymb}
\usepackage{amsmath}
\usepackage{braket}
\usepackage{enumerate}
\usepackage[margin=0.5in]{geometry}
\usepackage{xcolor}
\title{Introductory Quantum Mechanics Homework 3}
\author{Jeremy Chen}
\date{January 31, 2026}
\begin{document}
	\maketitle

\begin{enumerate}[\textbf{Problem \arabic{enumi}}]

\item Let \textbf{n} be a unit vector with polar coordinates $(\theta, \varphi)$ and let $\sigma$ denote the Pauli matrices. Find the eigenvectors of $\text{\textbf{n}} \cdot \sigma$. Hence show that the state of a spin-$\frac{1}{2}$ particle for which a measurement of spin along \textbf{n} yields $\hbar/2$ with uncertainty is 
$$\ket{\textbf{n}} = \sin\left( \frac{\theta}{2} \right) e^{i\varphi/2}\ket{\downarrow} + \cos\left( \frac{\theta}{2} \right)e^{-i\varphi/2}\ket{\uparrow}$$
where $\ket{\uparrow}$ and $\ket{\downarrow}$ are the usual eigenstates of $\hat{S}_{z}$. Obtain the corresponding expression for the eigenstate of $\textbf{n} \cdot \sigma$ with eigenvalue $-\hbar/2$. Explain why $\braket{\uparrow | \textbf{n}} = 0$ at $\theta = \pi$ and why both coefficients have modulus $1/\sqrt{2}$ when $\theta = \pi/2$.  

\color{violet}
The setup for $\textbf{n} \cdot \sigma = \sum_{i = 1}^{n} \textbf{n}_{i}\sigma_{i}$.
\begin{gather*}
\textbf{n} \cdot \sigma = \sin(\theta)\cos(\varphi)\sigma_{x} + \sin(\theta)\sin(\varphi)\sigma_{y} + \cos(\theta)\sigma_{z} \\
= \begin{bmatrix}
0 & \sin(\theta)\cos(\varphi) \\
\sin(\theta)\cos(\varphi) & 0
\end{bmatrix} + \begin{bmatrix}
0 & -i(\sin(\theta)\sin(\varphi)) \\
i(\sin(\theta)\sin(\varphi)) & 0
\end{bmatrix} + \begin{bmatrix}
\cos(\theta) & 0 \\
0 & -\cos(\theta)
\end{bmatrix} \\
= \begin{bmatrix}
\cos(\theta) & \sin(\theta)\cos(\varphi) - i\sin(\theta)\sin(\varphi) \\
\sin(\theta)\cos(\varphi) + i\sin(\theta)\sin(\varphi) & -\cos(\theta)
\end{bmatrix} \\
\begin{bmatrix}
\cos(\theta) & \sin(\theta)(\cos(\varphi) - i\sin(\varphi)) \\
\sin(\theta)(\cos(\varphi) + i\sin(\varphi)) & -\cos(\theta)
\end{bmatrix} \\
= \begin{bmatrix}
\cos(\theta) & \sin(\theta)e^{-i\varphi} \\
\sin(\theta)e^{i\varphi} & -\cos(\theta)
\end{bmatrix}
\end{gather*}
All of the Pauli matrices have eigenvalues $\pm 1$, so the eigenvectors would have the following forms. 
\begin{gather*}
\lambda = 1 \longrightarrow \begin{bmatrix}
	e^{-i\varpi}(\cot(\theta) + \csc(\theta)) \\
	1
\end{bmatrix} = \begin{bmatrix}
	e^{-i\varphi}\frac{\cos\left(\frac{\theta}{2}\right)}{\sin\left(\frac{\theta}{2}\right)} \\
	1
\end{bmatrix} \\
\lambda = -1 \longrightarrow \begin{bmatrix}
e^{-i\varphi}(\cot(\theta) - \csc(\theta))\\
1
\end{bmatrix} = \begin{bmatrix}
-e^{-i\varphi}\frac{\sin\left( \frac{\theta}{2} \right)}{\cos\left(\frac{\theta}{2}\right)}\\
1
\end{bmatrix}
\end{gather*}

\color{black}

\item For the harmonic oscillator calculate the energy eigenfunctions $\psi_{3}(x) = \braket{x | 3}$ and $\psi_{4}(x) = \braket{x | 4}$ explicitly. 

\color{violet}
The general function for solving some wavefunction at energy level $n$ is $\psi_{n}(x) = \frac{1}{\sqrt{2^{n}}n!}(y - \frac{d}{dy})^{n}\psi_{0}(y)$ where $y = \sqrt{\frac{m\omega}{\hbar}}x$. At the the third excited state, the following is the energy eigenfunction
\begin{gather*}
\psi_{3}(x) = \frac{1}{\sqrt{2^{3}3!}}\left(y - \frac{d}{dy}\right)^{3}\psi_{0}(y) = \frac{1}{\sqrt{48}}\left( \sqrt{\frac{m\omega}{\hbar}x} - \frac{d}{dy}  \right)^{3}\psi_{0}\left(\sqrt{\frac{m\omega}{\hbar}}x\right) \\
= \frac{1}{48}\left( \sqrt{\frac{m\omega}{\hbar}}x - \frac{d}{dy} \right)^{3} \left( \frac{m\omega}{\pi\hbar} \right)^{\frac{1}{4}} e^{-\frac{m\omega}{2\hbar}\frac{m\omega}{\hbar}x^{2}} = \frac{1}{48}\left( \sqrt{\frac{m\omega}{\hbar}}x - \frac{d}{dy} \right)^{3} \left( \frac{m\omega}{\pi\hbar} \right)^{\frac{1}{4}} e^{-\frac{m^{2}\omega^{2}}{2\hbar^{2}}x^{2}} 
\end{gather*}
At the fourth excited state, the following is the energy eigenfunction
\begin{gather*}
\psi_{4}(x) = \frac{1}{\sqrt{2^{4}4!}} \left( \sqrt{\frac{m\omega}{\hbar}}x - \frac{d}{dy} \right)^{4}\left( \frac{m\omega}{\pi\hbar} \right)^{\frac{1}{4}} e^{-\frac{m^{2}\omega^{2}}{2\hbar^{2}}x^{2}}
\end{gather*}
\color{black}

\item Consider a 1D quantum harmonic oscillator with classical frequency $\omega$. Define the coherent state $\ket{a}$ by 
$$\ket{\alpha} = e^{\alpha\hat{a}^{\dagger} - \alpha^{*}\hat{a} }\ket{0}$$
where $\hat{a}^{\dagger}$, $\hat{a}$ are the raising and lowering operators, $\ket{0}$ is the ground state, and $\alpha \in \mathbb{C}$. 

\begin{enumerate}

\item Show that $\ket{\alpha}$ is an eigenstate of $\hat{a}$ and find its eigenvalue. Compute the inner product between two different coherent states $\ket{\alpha}$ and $\ket{\beta}$. Does the set of all states $\ket{\alpha}$ for all $a \in \mathbb{C}$ form a basis of the Hilbert space?

\color{violet}
The current form of the coherent state could be reorganized in the following manner, such that $\ket{a}$ would be normalized. 
$$\ket{\alpha} = e^{\alpha\hat{a}^{\dagger} - \alpha^{*}\hat{a}}\ket{0} = e^{-\frac{1}{2}|\alpha|^{2}}e^{\alpha\hat{a}^{\dagger}}e^{-a^{*}\hat{a}}\ket{0} = e^{-\frac{1}{2}|\alpha|^{2}}\sum_{n = 0}^{\infty} \frac{\alpha^{n}}{\sqrt{n!}} \ket{n} = e^{-\frac{1}{2}|\alpha|^{2}}\sum_{n = 0}^{\infty} \frac{\alpha^{n}}{\sqrt{n!}}$$
By adding $\hat{a}$, the following is the new state.
$$\hat{a}\ket{\alpha} = e^{-\frac{1}{2}|\alpha|^{2}} \sum_{n = 0}^{\infty} \frac{\alpha^{n}}{\sqrt{n!}} \hat{a}\ket{n} = e^{-\frac{1}{2}|\alpha|^{2}} \sum_{n = 1}^{\infty} \frac{\alpha^{n}}{\sqrt{(n - 1)!}} \ket{n - 1} = \alpha\ket{\alpha}$$
The eigenvalue of $\ket{\alpha}$ is $\alpha$. Since $\ket{\alpha}$ could be represented in its normalized form, so could $\ket{\beta}$. The inner product of $\ket{\alpha}$ and $\ket{\beta}$ is the following
\begin{gather*}
\braket{\beta | \alpha} = e^{-\frac{1}{2}|\alpha|^{2} - \frac{1}{2}|\beta|^{2}}\braket{0 | e^{\beta^{*}\hat{a}}e^{\alpha\hat{a}^{\dagger}} | 0} = e^{-\frac{1}{2}|\alpha|^{2} - \frac{1}{2}|\beta|^{2}}\braket{0 | e^{\beta^{*}\hat{a} + \alpha\hat{a}^{\dagger} + \frac{1}{2}\beta^{*}\alpha} | 0} \\
= e^{-\frac{1}{2}|\alpha|^{2}-\frac{1}{2}|\beta|^{2} + \beta^{*}\alpha}\braket{0 | e^{\alpha\hat{a}^{\dagger}}e^{\beta^{*}\hat{a}} | 0} = e^{-\frac{1}{2}|\alpha|^{2}-\frac{1}{2}|\beta|^{2} + \beta^{*}\alpha}
\end{gather*}
Since $\alpha$ is the average position and momentum of the coherent state, it forms a basis of the Hilbert space. 
\color{black}

\item A quantum state is prepared in state $\ket{\alpha}$ at $t = 0$. Show that it evolves to become a new coherent state $e^{-i\omega t}\ket{e^{-i\omega t}\alpha}$. 

\color{violet}

\color{black}

\item Suppose $\alpha \in \mathbb{R}$. By expressing $\hat{a}$ and $\hat{a}^{\dagger}$ in terms of $\hat{x}$ and $\hat{p}$, sketch the position space wavefunction of $\ket{\alpha}$. 

\color{violet}
In terms of $\hat{x}$ and $\hat{p}$, $\hat{a}$ and $\hat{a}^{\dagger}$ are expressed like the following
$$\hat{a} = \sqrt{\frac{m\omega}{2\hbar}}\left( \hat{x} + \frac{i}{m\omega}\hat{p} \right),\ \ \hat{a}^{\dagger} = \sqrt{\frac{m\omega}{2\hbar}}\left( \hat{x} - \frac{i}{m\omega}\hat{p} \right)$$
These values of $\hat{a}$ and $\hat{a}^{\dagger}$ can be plugged into $\ket{\alpha}$ to get the following
\begin{gather*}
\ket{\alpha} = e^{-\frac{1}{2}|\alpha|^{2}}e^{\alpha\hat{a}^{\dagger}}\ket{0} = e^{-\frac{1}{2}|\alpha|^{2}}e^{\alpha \sqrt{\frac{m\omega}{2\hbar}}\left( \hat{x} - \frac{i}{m\omega}\hat{p} \right)}\ket{0} = e^{-\frac{1}{2}|\alpha|^{2}}e^{\alpha\frac{m\omega}{2\hbar}\hat{x} - \alpha\frac{m\omega}{2\hbar}\frac{i}{m\omega}\hat{p}}\ket{0}
\end{gather*}
\color{black}

\item Let $\alpha \in \mathbb{C}$ and compute $\braket{\alpha | \hat{p} | \alpha}$. Give a physical interpretation of the coherent state for general complex $\alpha$. Without futher calculation, describe the shape and motion of both the position space and momentum space wavefunctions as time passes. 

\color{violet}
The position wavefunction of a coherent state is defined as some $\psi_{\alpha}(x) = \braket{x | a}$. The expectation value of $x$ for the wavefunction would be $\int \psi_{\alpha}^{*} x \psi_{\alpha}dx$, which could be written in terms of ladder opperators as $\braket{x} = \braket{\alpha | \sqrt{\frac{2\hbar}{m\omega}}(\hat{a} + \hat{a}^{\dagger}) | \alpha}$. Momentum could be expressed in the same way, but instead $\braket{p} = \braket{\alpha | -i\sqrt{\frac{m\omega\hbar}{2}} (\hat{a} - \hat{a}^{\dagger})| \alpha}$. The it could be calculated like the following
$$\braket{\hat{p}} = -i\sqrt{\frac{m\omega\hbar}{2}}\braket{\alpha | (\hat{a} - \hat{a}^{\dagger}) | \alpha} = -i\sqrt{\frac{m\omega\hbar}{2}} (\braket{\alpha | \hat{a} | \alpha} - \braket{\alpha | \hat{a}^{\dagger} | \alpha}) = \sqrt{-\frac{m\omega\hbar}{2}} (\alpha - \alpha^{*}) = \sqrt{2m\omega\hbar}\text{Im}(\alpha)$$
The shape of the position and momentum wavefunctions remains the same as time passes, but they will begin to move towards the right of the graph. 
\color{black}

\end{enumerate}

\end{enumerate}

\end{document}